% ===== PHỤ LỤC A: BẢNG PHÂN CÔNG CÔNG VIỆC =====
\clearpage
\pagestyle{appendixlandscape}
\markboth{}{}

% -------- Trang 1: Tổng quan + Bảng A.1 --------
\begin{landscape}
\chapter{BẢNG PHÂN CÔNG CÔNG VIỆC}

Để bảo đảm tính minh bạch, công bằng và dễ dàng theo dõi tiến độ trong suốt quá trình thực hiện đồ án,
nhóm đã tiến hành phân chia nhiệm vụ một cách rõ ràng ngay từ giai đoạn đầu. Việc phân công cụ thể
giúp xác định trách nhiệm của từng thành viên, giảm thiểu sự chồng chéo trong công việc và đảm bảo
tất cả hoạt động đều được quản lý có hệ thống.

\begin{table}[htbp]
    \caption{Bảng tổng quan phân chia công việc trong nhóm}
    \centering
    \begin{tabularx}{\linewidth}{@{} c c X X c c c @{}}
    \toprule
    \textbf{STT} & \textbf{MSSV} & \textbf{Người thực hiện} & \textbf{Tóm tắt vai trò} &
    \textbf{Số lượng NV} & \textbf{Số lượng hoàn thành} & \textbf{Ghi chú} \\
    \midrule
    1 & [MSSV 1] & [Họ và tên Sinh viên 1] & Điều phối tiến độ nhóm, [Mô tả nhiệm vụ 1] & 5 & 5 & Hoàn thành \\[4pt]
    2 & [MSSV 2] & [Họ và tên Sinh viên 2] & [Mô tả nhiệm vụ 2] & 5 & 5 & Hoàn thành \\[4pt]
    3 & [MSSV 3] & [Họ và tên Sinh viên 3] & [Mô tả nhiệm vụ 3] & 5 & 5 & Hoàn thành \\[4pt]
    4 & [MSSV 4] & [Họ và tên Sinh viên 4] & [Mô tả nhiệm vụ 4] & 5 & 5 & Hoàn thành \\[4pt]
    \bottomrule
    \end{tabularx}
\end{table}
\end{landscape}

\clearpage
% \pagestyle{customfancy}   % phục hồi header/footer cho phụ lục B, C
